\documentclass[11pt]{amsart}

\usepackage{amsmath, amssymb, amsthm}
\usepackage{mathtools}
\usepackage{thmtools}
\usepackage{hyperref}

% % Theorem environments
\declaretheorem[numberwithin=section]
{theorem}
\declaretheorem{lemma, proposition, corollary}[
style=plain,
sibling=theorem,
]
\declaretheorem[style=definition]
{definition, example, condition}
\declaretheorem[style=remark]{remark, note, observation}


\title{Review on \\ McPhail--Snyder's work}
\author{knottedtail}
\address{Chonnam University}
\email{philip94@snu.ac.kr}
\date{13 April, 2026}


\begin{document}

\begin{abstract}
    Here, we review all the paper written by McPhail-Snyder.
\end{abstract}

\maketitle

\section{Timelines}

There are 8 papers written by McPhail--Snyder.
The following is citation of those papers in the order of time.
\begin{itemize}
    \item \cite{McPhailSnyder_Miller-2020-palnar_diagrams}
    \item \cite{Chen_McPhailSnyder_Morrison_Snyder-2021-Kashaev_Reshetikhin}
    \item \cite{McPhailSnyder-2021-thesis}
    \item \cite{McPhailSnyder-2022-holonomy_torsion}
    \item \cite{McPhailSnyder-2025-octahedral_coordinates}
    \item \cite{McPhailSnyder-2025-quantization_tangle}
    \item \cite{McPhailSnyder-2026-hyperbolic_octa_quantum}
    \item \cite{McPhailSnyder_Reshetikhin-2025-cyclic_modules}
\end{itemize}

In this paper, we review part of those papers related to the octahedral decomposition of knot complemnts and quantum invariants (holonomy invariants).

\section{\cite{McPhailSnyder_Miller-2020-palnar_diagrams}}
\section{\cite{Chen_McPhailSnyder_Morrison_Snyder-2021-Kashaev_Reshetikhin}}
\section{\cite{McPhailSnyder-2021-thesis}}
\section{\cite{McPhailSnyder-2022-holonomy_torsion}}
\section{\cite{McPhailSnyder-2025-octahedral_coordinates}}
\section{\cite{McPhailSnyder-2025-quantization_tangle}}
\section{\cite{McPhailSnyder-2026-hyperbolic_octa_quantum}}
\section{\cite{McPhailSnyder_Reshetikhin-2025-cyclic_modules}}


\bibliographystyle{amsalpha}
% \nocite{*}
\bibliography{bibliography}

\end{document}
